\documentclass[10pt,a4paper]{amsart}
\usepackage[margin=19mm]{geometry}
\usepackage{amsmath,amssymb,mathtools}
\usepackage[scr=rsfs,cal=euler]{mathalfa}
\usepackage{xfrac}
\usepackage[unicode,hidelinks,bookmarks=false]{hyperref}
\newcommand{\ie}{i.e.\ }
\newcommand{\cf}{cf.\ }
\newcommand{\resp}{respectively\ }
\newcommand{\Subsection}{\subsection}
\providecommand{\nobreakdash}{\nobreak\hbox{-}}
\setlength{\emergencystretch}{2em}
\multlinegap0pt
\usepackage{bm}
\usepackage{bbm}
\usepackage{dsfont}
\usepackage{xspace}
\usepackage{cancel}
\usepackage{mathtools}
\usepackage[shortlabels]{enumitem}
\usepackage{amsthm}
\makeatletter
\@ifundefined{claim}{\newtheorem{claim}{Claim}}{}
\@ifundefined{theorem}{\newtheorem{theorem}{Theorem}}{}
\makeatother
\title[Some questions on entangled linear orders]{Some questions on entangled linear orders}
\author{Rapha\"el Carroy, Maxwell Levine, Lorenzo Notaro}
\date{Source: arXiv:2507.17503v2 (CC BY 4.0)}
\begin{document}
\maketitle

\noindent\emph{Source and licence.} Some questions on entangled linear orders. Version arXiv:2507.17503v2 (CC BY 4.0), distributed under the Creative Commons Attribution 4.0 International licence, as recorded in the arXiv licence metadata for this exact version and shown on the abstract page https://arxiv.org/abs/2507.17503v2. Full rights notice: licenses/arxiv-2507.17503-CC-BY-4.0.txt.\par\smallskip

\noindent\textit{Editorial scope.} Counted unit: the Theorem whose source label is \texttt{None} in the article (source0.tex lines 704--706, source label \texttt{None}), delivered together with the article's own complete original proof of it (source0.tex lines 707--870, one \texttt{proof} environment of 164 lines). No mathematics is written, repaired or re-derived: statement and proof are the article's own text line for line. Required context retained uncounted: none. Every definition, display and label of those lines is retained unchanged. Omitted from this package and not redistributed: 1--703, 871--880. Nothing inside a retained excerpt is omitted or altered: every retained body is delivered exactly as the article's lines are written, so no omission receipt of any kind accompanies this package. Citations: the retained excerpts carry no citation command at all, so no bibliography entry of the article is transcribed and no reference is redistributed. Reference adaptation: none was needed and none was made; every reference inside the delivered unit resolves inside it, so no unresolved reference or citation is produced. Printed slips kept literal: no printed word, symbol, hypothesis or number of the counted statement or of its proof was altered. Layout and class: the article's own class is \texttt{amsart} with packages including graphicx, amsmath, amsthm, amsfonts, amssymb, enumitem, mathtools, biblatex, hyperref, cleveref, fixme; the wrapper is the lane's pinned \texttt{amsart} editorial wrapper at 10pt on A4, which restates the theorem-like environments the retained text needs and adds the article's own macro definitions that the retained text needs (none), with extra wrapper packages (bm, bbm, dsfont, xspace, cancel, mathtools, enumitem). The delivered numbering is the unit's own. Assets: no retained span contains a figure environment, a graphics-inclusion command, a PDF-page inclusion, a table or a TikZ picture; no image file is copied into this package, the package declares no asset dependency and no image file is redistributed with it. Licence: version arXiv:2507.17503v2 of this article is distributed under the Creative Commons Attribution 4.0 International licence, as recorded in the arXiv licence metadata for this exact version and shown on the abstract page https://arxiv.org/abs/2507.17503v2. A full rights notice is delivered at licenses/arxiv-2507.17503-CC-BY-4.0.txt. Characters: every retained excerpt is pure ASCII, so the delivered engine, XeLaTeX with its default Latin Modern text font, reports no missing character.
\begin{theorem}
Suppose that $\diamondsuit$ holds. Then, there exists a $2$-entangled non-separable linear order.
\end{theorem}

\begin{proof}
The construction is similar to the classical construction of a Suslin line under $\diamondsuit$. We define a tree ordering $\unlhd$ on $\omega_1\times \omega$ with $\mathrm{Lev}_\alpha(\omega_1 \times \omega, \unlhd) = \{\alpha\}\times \omega$ for every $\alpha$. For every $x = (\alpha, k) \in \omega_1\times\omega$, we denote $\alpha$ (resp. $k$) by $\mathrm{ht}(x)$ (resp. $n(x)$). Moreover, given $x,y \in \omega_1\times \omega$, we write $x \top y$ (resp. $x \perp y$) when $x$ and $y$ are comparable (resp. incomparable) with respect to $\unlhd$.

We want to define the tree $(\omega_1\times \omega, \unlhd)$ such that the following lexicographic-like linearization of $\unlhd$ is $2$-entangled and non-separable: for all $x,y \in \omega_1\times \omega$,
\begin{multline*}
x \preceq y \text{ if either } x \unlhd y \text{ or } \big( x \perp y \text{ and } \\
n(\min\{z \unlhd x \mid z \not\trianglelefteq y\}) >
n(\min \{z \unlhd y \mid z \not\trianglelefteq x\})\big),
\end{multline*}

\smallskip
\noindent where the minimum is taken with respect to $\unlhd$. From now on, the minima and maxima between nodes of our tree are taken with respect to $\preceq$.

We now introduce a useful binary relation induced by $\preceq$: given $x,y \in \omega_1 \times \omega$, we let $x \prec^+ y$ if $x \preceq y$ and $x \perp y$. The following straightforward observations are often implicitly used afterward: if $x \prec^+ y$, then $z \prec^+ y$ for all $z \unrhd x$; moreover, if $y \preceq x$, then $y \preceq z$ for all $z \unrhd x$.

For the sake of clarity, we denote by $L_\alpha$ the set  $\{\alpha\}\times \omega$ for each $\alpha$. Fix a $\diamondsuit$-sequence $\langle A_\alpha \mid \alpha < \omega_1\rangle$. Fix also a bijection $\psi: \omega_1 \rightarrow (\omega_1 \times \omega)^2$.  We define $\unlhd$ inductively such that:\vspace{0.3em}
\begin{enumerate}
\itemsep0.3em
\item $\unlhd$ is a tree ordering on $\omega_1 \times \omega$ and $\mathrm{Lev}_\alpha(\omega_1 \times \omega, \unlhd) = L_\alpha$ for every $\alpha < \omega_1$,
\item for every $n$ and $\alpha < \omega_1$, there are infinitely many $m$'s such that $(\alpha, n) \unlhd (\alpha+1, m)$,
\item if $\beta < \alpha < \omega_1$ and $x \in L_\beta$, then there is a $y \in L_\alpha$ with $x \lhd y$,
\item Suppose that $\alpha$ is limit and $\psi``A_\alpha \subseteq (\alpha\times \omega)^2$ is a map, denoted for clarity by $f_\alpha$, which is maximal among the 1-1 and monotone maps $f$ from a subset of $\alpha \times \omega$ to $\alpha \times \omega$ such that $f(x) \neq x$ for all $x \in \mathrm{dom}(f)$. Then, for all $x,y \in (\alpha+1) \times \omega$, if either

\smallskip
\begin{enumerate}[label=(4\alph*)]
\item $\mathrm{ht}(x) = \mathrm{ht}(y) = \alpha$, or
\item $\mathrm{ht}(x) < \alpha$ and $\mathrm{ht}(y) = \alpha$ and $x \not\in \mathrm{dom}(f_\alpha)$, or
\item $\mathrm{ht}(x) = \alpha$ and $\mathrm{ht}(y) < \alpha$ and $y \not\in \mathrm{ran}(f_\alpha)$,
\end{enumerate}

\smallskip
\noindent then the map $f_\alpha \cup \{(x,y)\}$ is not monotone.\vspace{0.3em}
\end{enumerate}

In (4) and onward, monotonicity is always with respect to $\preceq$.
Suppose for the moment that we have already defined $\unlhd$ satisfying (1)-(4). We claim that the linear order $(\omega_1 \times \omega, \preceq)$ is $2$-entangled and non-separable.

\begin{claim}
$(\omega_1 \times \omega, \preceq)$ is non-separable.
\end{claim}
\begin{proof}
Fix a countable $D\subseteq \omega_1 \times \omega$. Then, there is an $\alpha < \omega_1$ such that $D\subseteq \alpha \times \omega$. By (2), we can pick $x \in L_\alpha$  and $y \in L_{\alpha+1}$  with $x \lhd y$. Then, the open interval $]x, y[$ in $(\omega_1 \times \omega, \preceq)$ is nonempty as, by (2) again, it contains infinitely many elements of $L_{\alpha+1}$. Moreover, it is straightforward that ${]x,y[} \subseteq (\omega_1\setminus \alpha) \times \omega$. Thus, no element of $D$ belongs to $]x,y[$. We conclude that $(\omega_1 \times \omega, \preceq)$ is non-separable.
\end{proof}


\begin{claim}
$(\omega_1 \times \omega, \preceq)$ is $2$-entangled.
\end{claim}
\begin{proof}
Fix a 1-1 monotone map $f: A \rightarrow \omega_1 \times \omega$, with $A \subseteq \omega_1 \times \omega$, such that $f(x) \neq x$ for all $x \in A$ towards showing that $A$ is countable.

Let $\tilde{f}$ be a maximal 1-1 and monotone extension of $f$ such that $\tilde{f}(x) \neq x$ for all $x \in \mathrm{dom}(\tilde{f})$.


There is a club $C\subseteq\omega_1$ such that, for all $\alpha \in C$, $\alpha \times \omega$ is closed under both $\tilde{f}$ and $\tilde{f}^{-1}$ and  $\tilde{f}  \upharpoonright \alpha \times \omega$ is maximal among the maps from a subset of $\alpha \times \omega$ to $\alpha\times \omega$ that satisfy the same relevant properties of $\tilde{f}$---that is, being 1-1, monotone and pointwise different from the identity. By the properties of the $\diamondsuit$-sequence there is an $\alpha \in C$ with $f_\alpha = \tilde{f} \upharpoonright \alpha \times \omega$. We claim that $\mathrm{dom}(\tilde{f}) \subseteq \alpha \times \omega$, from which the countability of $A$ directly follows.

Suppose towards a contradiction that there exists some $x \in \mathrm{dom}(\tilde{f}) \setminus (\alpha \times \omega)$. Fix one such $x$. Note that, since $\alpha \times \omega$ is closed under $\tilde{f}^{-1}$, we also have $\tilde{f}(x) \not\in \alpha \times \omega$.  Let $x \upharpoonright \alpha$ be the $\unlhd$-predecessor of $x$ of height $\alpha$, and define $\tilde{f}(x) \upharpoonright \alpha$ analogously for $\tilde{f}(x)$.  By (4), the map $(\tilde{f} \upharpoonright \alpha \times \omega) \cup \{(x \upharpoonright \alpha, \tilde{f}(x) \upharpoonright \alpha)\}$ is not monotone.

It is straightforward that for every $z,w \in \omega_1\times \omega$ with $\mathrm{ht}(z) > \mathrm{ht}(w)$, $z \prec w$ if and only if $z \prec^+ w$. Therefore, it follows directly from the useful observation made after the definition of $\prec^+$ and from $(\tilde{f} \upharpoonright \alpha \times \omega) \cup \{(x \upharpoonright \alpha, \tilde{f}(x) \upharpoonright \alpha)\}$ not being monotone that also $(\tilde{f} \upharpoonright \alpha \times \omega) \cup \{(x, \tilde{f}(x))\}$ is not monotone, which is a contradiction.
\end{proof}

We define $\unlhd$ by induction on its levels. Suppose that we have defined $\unlhd$ on $\alpha \times \omega$ so that (1)-(4) holds below $\alpha$, towards extending $\unlhd$ to $(\alpha+1) \times \omega$.

We deal with the only nontrivial case: $\alpha$ is limit, and the hypotheses of (4) are satisfied. Fix an increasing sequence $(\xi_m)_{m\in\omega}$ of countable ordinals such that $\sup_m \xi_m = \alpha$ and fix also an enumeration $(b_n)_{n\in\omega}$ of $\alpha \times \omega$.

We define a family $\langle z^n_m \mid n,m\in\omega\rangle$ of elements of $\alpha \times \omega$ such that $\mathrm{ht}(z^n_m) \ge \xi_m$ and $z^n_{m} \unlhd z^n_{m+1}$ for every $n,m\in\omega$. Once we have defined this family, we extend $\unlhd$ as follows: for every $x \in \alpha \times \omega$ and $n\in\omega$, we let $x \lhd (\alpha, n)$ if $x \unlhd z^n_m$ for some $m$.

We describe how to define the $z^{n}_m$s and only then we show that this definition works---that is, the resulting extension of $\unlhd$ on $(\alpha+1) \times \omega$ still satisfies (1)-(4) below $\alpha+1$.

Fix an $n\in\omega$ and suppose that we have defined $z^k_m$ for all $k < n$ and all $m$, towards defining $z^{n}_m$ for all $m$. Fix an enumeration $(w_m, i_m)_{m\in\omega}$ of the set $((\alpha \times \omega) \cup (\{\alpha\} \times n))\times 2$.


First let us define $z_0^n$. Suppose first that there is no $z \in \mathrm{dom}(f_\alpha) \cup \mathrm{ran}(f_\alpha)$ such that  $b_n \lhd z$: in this case let $z_0^n$ be such that $b_n \lhd z_0^n$ and $\mathrm{ht}(z_0^n) = \max(\mathrm{ht}(b_n), \xi_0)+2$---we can find such $z_0^n$ by (3). 

Now suppose that there is a $z \in \mathrm{dom}(f_\alpha) \cup \mathrm{ran}(f_\alpha)$ such that  $b_n \lhd z$. Fix one such $z$. Let $z^\star = f_\alpha(z)$ if $z \in \mathrm{dom}(f_\alpha)$ and $z^\star = f_{\alpha}^{-1}(z)$ otherwise.  We need to consider the following cases:

%We want to find some $z_0^n \rhd b_n$ with $\mathrm{ht}(z_0^n) \ge \xi_0$ such that, for some $x \in \mathrm{dom}(f_\alpha)$, if $f_\alpha$ is increasing (resp. decreasing), the pairs $(z_0^n, z_0^n)$ and $(x, f_\alpha(x))$ realize the type $\langle 1, 0\rangle$ (resp. $\langle 0,0\rangle$) in the following strong sense: either $z_0^n \prec^+ x$ and $f_\alpha(x) \preceq z_0^n$ (resp. $z_0^n \prec^+ f_\alpha(x)$), or $x \preceq z_0^n$ and $z_0^n \prec^+ f_\alpha(x)$ (resp. $f_\alpha(x) \preceq z_0^n$). In order to do so, we need to consider different cases:

\begin{enumerate}[label={(\roman*)}]
\itemsep0.3em
\item $z \top z^\star$ and $f_\alpha$ is increasing:
\newline by the way we chose $z$ and by case assumption, $z,z^\star$ and $b_n$ are all $\unlhd$-compatible and, moreover, $\max(\min(z,z^\star), b_n) \lhd \max(z,z^\star)$. Let $\delta = \mathrm{ht}(\max(\min(z, z^\star), b_n))+1$, and let $i \in \omega$ be such that $(\delta, i) \unlhd \max(z, z^\star)$. By (2), we can pick $j > i$ such that $\max(\min(z, z^\star), b_n) \lhd (\delta, j)$. By (3) we can pick $z_0^n$ such that $(\delta, j) \unlhd z_0^n$ and $\mathrm{ht}(z_0^n) \ge \xi_0$. In particular, we have $\min(z,z^\star) \prec z_0^n \prec^+ \max(z, z^\star)$. 

\item $z \top z^\star$ and $f_\alpha$ is decreasing:
\newline by (3) we can pick $z_0^n$ such that $\max (z,z^\star) \unlhd z_0^n$ and $\mathrm{ht}(z_0^n) \ge \xi_0$. In particular, we have $z,z^\star \preceq z_0^n$.

\item $z \prec^+ z^\star$ and $b_n \not\trianglelefteq z^\star$ and $f_\alpha$ is increasing:
\newline by (3) we can pick $z_0^n$ such that $z \unlhd z_0^n$ and $\mathrm{ht}(z_0^n) \ge \xi_0$. In particular, we have $z \preceq z_0^n \prec^+ z^\star$.


\item $z^\star \prec^+ z$ and $b_n \not\trianglelefteq z^\star$ and $f_\alpha$ is increasing:
\newline let $i \in \omega$ be such that $(\mathrm{ht}(b_n)+1, i) \unlhd z$. By (2), we can fix some $j > i$ such that $b_n \lhd (\mathrm{ht}(b_n)+1, j)$. By (3), we can pick $z_0^n$ such that $(\mathrm{ht}(b_n)+1, j) \unlhd z_0^n$ and $\mathrm{ht}(z_0^n) \ge \xi_0$. In particular, we have $z^\star \prec z_0^n \prec^+ z$.


\item $z \perp z^\star$ and $b_n \unlhd z, z^\star$ and $f_\alpha$ is increasing:
\newline by (3), we can pick $z_0^n$ such that $\min (z,z^\star) \unlhd z_0^n$ and $\mathrm{ht}(z_0^n) \ge \xi_0$. In particular, we have $\min (z,z^\star) \preceq z_0^n \prec^+ \max (z,z^\star)$.


\item $z \prec^+ z^\star$ and $b_n \not\trianglelefteq z^\star$ and $f_\alpha$ is decreasing:
\newline pick $z_0^n$ as in (iv). In particular, we have $z_0^n \prec^+ z,z^\star$.


\item $z^\star \prec^+ z$ and $b_n \not\trianglelefteq z^\star$ and $f_\alpha$ is decreasing:
\newline pick $z_0^n$ as in (iii). In particular, we have $z,z^\star \preceq z_0^n$.


\item $z \perp z^\star$ and $b_n \unlhd z, z^\star$ and $f_\alpha$ is decreasing:
\newline by (3), we can pick $z_0^n$ such that $\max (z,z^\star) \unlhd z_0^n$ and $\mathrm{ht}(z_0^n) \ge \xi_0$. In particular, we have $z,z^\star \preceq z_0^n$.
\end{enumerate}




Suppose that we have defined $z_m^n$ towards defining $z_{m+1}^n$. If either $i_m = 0$ and $w_m \in \mathrm{dom}(f_\alpha)$, or $i_m = 1$ and $w_m \in \mathrm{ran}(f_\alpha)$, then simply let $z_{m+1}^n$ be such that $z_{m}^n \unlhd z_{m+1}^n$ and $\mathrm{ht}(z_{m+1}^n) \ge \xi_{m+1}$---we can find such a $z_{m+1}^n$ by (3). 

Now assume that either $i_m = 0$ and $w_m \not\in \mathrm{dom}(f_\alpha)$, or $i_m = 1$ and $w_m \not\in \mathrm{ran}(f_\alpha)$. If there is no $z$ with $z_m^n \lhd z$ such that $z \in \mathrm{ran}(f_\alpha)$ if $i_m = 0$ and $z \in \mathrm{dom}(f_\alpha)$ if $i_m = 1$, then let $z_{m+1}^n$ be such that $z_m^n \lhd z_{m+1}^n$ and $\mathrm{ht}(z_{m+1}^n) \ge \xi_{m+1}$ and $z_{m+1}^n \neq w_m$. Otherwise, fix one such $z$ and let $z^\star$ denote $f^{-1}_\alpha(z)$ if $i_m  = 0$, and $f_\alpha(z)$ if $i_m = 1$. Now we need to consider the following cases:
\begin{enumerate}[label={(\alph*)}]
\itemsep0.3em
\item $f_\alpha$ is increasing and $z^\star \prec w_m$:
\newline let $i \in \omega$ be such that $(\mathrm{ht}(z_m^n)+1, i) \unlhd z$. By (2), there we can fix some $j > i$ such that $z_m^n \lhd (\mathrm{ht}(z_m^n)+1, j)$. By (3), we can pick $z_{m+1}^n$ such that $(\mathrm{ht}(z_m^n)+1, j) \unlhd z_{m+1}^n$ and $\mathrm{ht}(z_{m+1}^n) \ge \xi_{m+1}$. In particular, we have $z_{m+1}^n \prec^+ z$


\item $f_\alpha$ is decreasing and $w_m \prec z^\star$:
\newline as in (a).

\item $f_\alpha$ is increasing and $w_m \prec z^\star$:
\newline by (3), we can pick $z_{m+1}^n$ such that $z \unlhd z_{m+1}^n$ and $\mathrm{ht}(z_{m+1}^n) \ge \xi_{m+1}$.


\item $f_\alpha$ is decreasing and $z^\star \prec w_m$:
\newline as in (c).
\end{enumerate}


This ends the inductive definition of the $z_m^n$s. We are left to prove that (1)-(4) are still satisfied below $\alpha+1$. 

Since we chose the $z_m^n$s such that $\sup_m \mathrm{ht}(z_m^n) = \sup_m \xi_m = \alpha$ for every $n$, we conclude that $\mathrm{Lev}_\alpha((\alpha+1)\times \omega, \unlhd) = L_\alpha$---that is, (1) is satisfied. 

Condition (2) trivially holds since it holds below $\alpha$ by the induction hypothesis, and $\alpha$ is limit.

Moreover, since $b_n \unlhd z_0^n \lhd (\alpha, n)$ for every $n$ by choice of $z_0^n$, (3) also holds---that is, every element of $\alpha\times \omega$ has an $\unlhd$-successor in $L_\alpha$.  


Finally, we need to check that (4) is satisfied. We prove it by induction in the following sense: we suppose that for all $x,y \in (\alpha \times \omega) \cup (\{\alpha\} \times n)$ satisfying one among (4a)-(4c), the map $f_\alpha \cup \{(x,y)\}$ is not monotone, towards showing that the same can be said for all $x,y$ taken in $(\alpha \times \omega) \cup (\{\alpha\} \times (n+1))$.

We will check only some representative cases, as the other are analogous. From now on we assume that $f_\alpha$ is increasing.

Let us show that $f_\alpha \cup \{((\alpha, n), (\alpha, n))\}$ is not increasing.  Suppose that there is a $z \in \mathrm{dom}(f_\alpha) \cup \mathrm{ran}(f_\alpha)$ such that  $b_n \lhd z$. Fix the $z$ satisfying this property we chose when defining $z_0^n$---we assume for simplicity $z \in \mathrm{dom}(f_\alpha)$, as the case $z \in \mathrm{ran}(f_\alpha)$ is analogous. Then, looking at the possible cases (i.e. (i), and (iii)-(v)), either $z_0^n \prec^+ z$ and $f_\alpha(z) \preceq z_0^n$, or $z_0^n \prec^+ f_\alpha(z)$ and $z \preceq z_0^n$. In particular, since $z_0^n \lhd (\alpha, n)$, either $(\alpha, n) \prec z$ and $f_\alpha(z) \prec (\alpha, n)$, or $z \prec (\alpha, n)$ and $(\alpha, n) \prec f_\alpha(z)$. In other words, the pairs $((\alpha, n), (\alpha, n))$ and $(z, f_\alpha(z))$ realize the type $\langle 1,0\rangle$. 

Now suppose instead that there is no $z \in \mathrm{dom}(f_\alpha) \cup \mathrm{ran}(f_\alpha)$ such that $b_n \lhd z$. By definition of $z_0^n$ and case assumption, we have $\mathrm{ht}(z_0^n) \ge \mathrm{ht}(b_n)+2$. Therefore, we can fix some $t$ such that $b_n \lhd t \lhd z_0^n$. By case assumption, both $t$ and $z_0^n$ do not belong to $\mathrm{dom}(f_\alpha) \cup \mathrm{ran}(f_\alpha)$. By maximality of $f_\alpha$, we conclude that there is an $x \in \mathrm{dom}(f_\alpha)$ such that $(x, f_\alpha(x))$ and $(t, z_0^n)$ realize $\langle 1,0\rangle$.From the latter observation and our case assumption, it follows that either $t \prec^+ x$ and $f_\alpha(x) \preceq z_0^n$, or $x \preceq t$ and $z_0^n \prec^+ f_\alpha(x)$.
Since $t \lhd z_0^n \lhd (\alpha, n)$, we conclude that the pairs $((\alpha, n), (\alpha, n))$ and $(x, f_\alpha(x))$ realize the type $\langle 1,0\rangle$. 

We have shown that $f_\alpha \cup \{((\alpha, n), (\alpha, n))\}$ is not increasing. Now fix some $k$ with $k < n$ and let us show that $f_\alpha \cup \{(\alpha, k), (\alpha, n)\}$ is not increasing. Fix  an $m$ such that $w_m = (\alpha, k)$ and $i_m = 0$. Clearly $w_m = (\alpha, k) \not\in \mathrm{dom}(f_\alpha)$.

Suppose first that there is a $z$ with $z_m^n \lhd z$ such that $z \in \mathrm{ran}(f_\alpha)$. Fix the $z$ satisfying this property we chose when defining $z_{m+1}^n$. Looking at the cases (a) and (c) in the definition of $z_{m+1}^n$, it is easy to see that either $f_\alpha^{-1}(z) \prec (\alpha, k)$ and $z_{m+1}^n \prec^+ z$, or $(\alpha, k) \prec f_\alpha^{-1}(z)$ and $z \preceq z_{m+1}^n$. Hence, the pairs $((\alpha, k), (\alpha, n))$ and $(f_\alpha^{-1}(z), z)$ realize $\langle 1,0\rangle$. 

Now suppose instead that there is no $z$ with $z_m^n \lhd z$ such that $z \in \mathrm{ran}(f_\alpha)$. Clearly, this means that $z_{m+1}^n \not\in \mathrm{ran}(f_\alpha)$. By the induction hypothesis, there is some $x \in \mathrm{dom}(f_\alpha)$ such that the pairs $(x, f_\alpha(x))$ and $((\alpha, k), z_{m+1}^n)$ realize $\langle 1,0\rangle$. But now note that, again by case assumption, we must have either $z_{m+1}^n \prec^+ f_\alpha(x)$ or $f_\alpha(x) \prec z_{m+1}^n$. Therefore, the pairs $(x, f_\alpha(x))$ and $((\alpha, k), (\alpha, n))$ also realize $\langle 1,0\rangle$.

We have shown that $f_\alpha \cup \{((\alpha, k), (\alpha, n))\}$ is not increasing. Now fix some $y \in \alpha \times \omega$ with $y \not\in \mathrm{ran}(f_\alpha)$ towards showing that $f_\alpha \cup \{((\alpha, n), y)\}$ is not increasing. Fix an $m$ such that $w_m = y$ and $i_m = 1$. 

If there is a $z$ with $z_m^n \lhd z$ such that $z \in \mathrm{dom}(f_\alpha)$, the argument used to show that there is some $x \in \mathrm{dom}(f_\alpha)$ such that  $(x, f_\alpha(x))$ and $((\alpha, n), y)$ realize $\langle 1,0\rangle$ is the same as in the previous case. So, let us assume that there is no such $z$. In particular, we must have $z_{m+1}^n \not\in \mathrm{dom}(f_\alpha)$. Moreover, by the way we chose $z_{m+1}^n$, we also have $z_{m+1}^n \neq w_m = y$. By maximality of $f_\alpha$, there is some $x \in \mathrm{dom}(f_\alpha)$ such that the pairs $(x, f_\alpha(x))$ and $(z_{m+1}^n, y)$ realize $\langle 1,0\rangle$. But now note that, again by case assumption, we must have either $z_{m+1}^n \prec^+ x$ or $x \prec z_{m+1}^n$. Therefore, the pairs $(x, f_\alpha(x))$ and $((\alpha, n), y)$ also realize $\langle 1,0\rangle$.

The other cases are analogous.
\end{proof}

\end{document}
